\section{Comparative Interpretation}\label{sec:comparative-interpretation}

To fully evaluate the utility and positioning of Motivo, it must be contextualized within the existing
ecosystem of music programming tools.
This section provides a comparative analysis against two distinct paradigms: live-coding environments (exemplified by
Sonic Pi) and low-level scripting libraries (such as Python's \texttt{mido}).

\subsection{Comparison with Sonic Pi}\label{detail-subsec:comparison-with-sonic-pi}

Sonic Pi is a highly successful live-coding language embedded in Ruby, designed primarily for performance and education.
While both Sonic Pi and Motivo allow music to be expressed as text, their underlying design philosophies lead
to different compositional workflows.

\subsubsection{Execution Model}
Sonic Pi utilizes a real-time, interpreted execution model.
The musician writes scripts that trigger sounds ``on the fly'' through the SuperCollider synthesis engine.
Time is managed explicitly by the programmer using the \texttt{sleep} function, which halts the execution thread.

In contrast, Motivo is fully compiled and ``zero-runtime''.
There is no concept of a thread sleeping; instead, the \texttt{play} and \texttt{rest} commands define a mathematical
relationship on a static timeline.
While Sonic Pi excels in improvisational settings where the code is meant to be modified while it runs,
Motivo is oriented toward \textit{structural composition}, where the composer needs a complex arrangement to
resolve into a consistently quantized MIDI file every time it is compiled.

\subsubsection{Structural Abstraction in Practice}
In Sonic Pi, polyphony and parallel execution are achieved by spawning new threads (e.g., using \texttt{in\_thread}).
This can lead to synchronization issues if the composer does not meticulously manage the timing of each thread.
Motivo resolves parallelism through the \texttt{voice} construct.
Because voices are unrolled at compile-time by isolating the internal cursor, the compiler supports that parallel
musical lines are aligned according to the compiler cursor model at their inception without requiring manual
thread management or timing arithmetic from the composer.

\subsection{Comparison with Raw MIDI Scripting}\label{detail-subsec:comparison-with-raw-midi-scripting}

For offline composition, many developers use general-purpose programming languages combined with MIDI libraries, such as
Python's \texttt{mido} or C++'s \texttt{RtMidi}.
While these tools offer absolute control over the output, they operate at an exceptionally low level of abstraction.

\subsubsection{Event Placement and Musical Intent}
When using a library like \texttt{mido}, the composer must manually construct \texttt{Note-On} and \texttt{Note-Off}
messages.
Crucially, because the MIDI file format relies on delta-times, the composer is forced to manually calculate the tick
difference between every sequential event across all tracks.

Motivo abstracts this entirely.
A composer writing \texttt{play (C4, E4, G4):2;} in Motivo communicates their \textit{musical intent}---playing a
C-Major triad for two beats.
Motivo's lowering engine and backend automatically handle the translation into absolute beats, the flattening into
individual notes, the sorting by absolute tick, and the final variable-length delta-time encoding.

\subsubsection{Boilerplate and Expressiveness}
A simple four-beat loop of a hi-hat playing eighth notes requires a significant amount of boilerplate in Python:
creating a file, creating a track, appending tempo metadata, and writing a loop that interleaves \texttt{Note-On} and
\texttt{Note-Off} messages with explicit delta times.

Motivo achieves the same result with minimal syntax:
\begin{lstlisting}[language=Motivo,label={lst:midi-scripting-comparison}]
tempo 120;
track using drums {
    loop (8) { play hihat :0.5; }
}
\end{lstlisting}

This contrast in verbosity illustrates Motivo's intended role as a domain-specific tool.
By embedding musical semantics (like default durations, rest tracking, and automatic track/channel assignment) directly
into the language design, the compiler reduces cognitive friction and allows the user to operate strictly within the
musical domain.
